% ===================================================================
%  نقشہ نمبر 9 — پہاڑ تے پاݨیاں والا شہر۔ — جنگل تے شکار۔
%  Traduction 2026 d'apres texto/io, controlee sur texto/fr, texto/en
%  et texto/hi. Consigne : voir texto/pnb/10-naqsha-01.tex.
% ===================================================================

%%K t09-apar-1 apar
\VUsaut{4.0mm}
\VUtitre{15.0pt}{60}{نقشہ نمبر 9}
\VUsaut{3.0mm}

%%K t09-tit-1 sub
\VUcentre{12.0pt}{0}{\textit{پہلا منظر۔}}
\VUsaut{3.0mm}

%%K t09-tit-2 sub
\VUpk{11.4pt}{40}{پہاڑ۔ --- پاݨیاں والا شہر۔}
\VUsaut{3.0mm}

%%K t09-c1-01-1 p
1. --- میں اک ہفتے توں پراتز وچ آں۔ پرینیز دی حیران کر دیݨ والی
\VUgras{پہاڑی لڑی} \textsuperscript{(1)} دے سامݨے کھڑا ہو کے مینوں جیہڑی حیرت
ہوئی، اوہ میں دس نئیں سکدا؛ اوہدی \VUgras{چوٹیاں دی لکیر} \textsuperscript{(2)}
نیلے اسمان اُتے کِسے وڈی جھالر وانگوں اُبھردی اے۔

%%K t09-c1-02-1 p
2. --- میں سوچیا وی نہ سی پئی پرینیز دیاں \VUgras{سِکھراں} \textsuperscript{(3)}
اِنّی \VUgras{اُچائی} تیکر اپڑ سکدیاں نیں، تے میں اوہناں شاندار
\VUgras{برفانی دریاواں} \textsuperscript{(5)} تے برف نال ڈھکیاں \VUgras{چوٹیاں}
\textsuperscript{(4)} نوں ویکھ کے دنگ رہ گیا، جیہناں توں اِنّے بھیانک
\VUgras{برف دے تودے} اُتردے نیں۔

%%K t09-tit-3 sub
\VUsaut{3.0mm}
\VUpk{10.2pt}{40}{پہاڑی منظر۔}
\VUsaut{3.0mm}

%%K t09-c1-03-1 p
3. --- آوݨ دے اگلے ای دن میں پراتز کول اوس پرانے بُجھے \VUgras{آتش فشاں}
\textsuperscript{(6)} تیکر گیا۔ \VUgras{مونہہ} دے ننگے بنجر کنڈھے اُتے ہُݨ وی اوہ
نشان صاف وکھائی دیندے نیں جیہڑے اوس \VUgras{لاوے} نے چھڈے سن، جیہڑا کئی صدیاں
پہلاں \VUgras{پہاڑ دی ڈھال} \textsuperscript{(7)} توں وگیا سی۔ اک
\VUgras{ڈھلوان} اُتے مینوں اوس لاوے دے کجھ ٹوٹے مل ای گئے۔

%%K t09-c1-04-1 p
4. --- پراتز سرحد اُتے وسیا اے، تے جیہڑا \VUgras{قلعہ} \textsuperscript{(8)} فرانس
تے ہسپانیہ دے وچکارلے \VUgras{درے} \textsuperscript{(27)} دی رکھوالی کردا اے، اوہ
راین دے کنڈھے دے اوہناں پرانے قلعیاں جیہا ای اے جیہڑے اپݨی چٹان دے اُتوں کِسے
\VUgras{شکار} دی تاک وچ لگدے نیں۔

%%K t09-c1-05-1 p
5. --- اک وڈا \VUgras{عقاب} \textsuperscript{(9)}، جیدا آہلݨا \VUgras{گڑھی}
\textsuperscript{(8)} توں دور نئیں، اکثر میرا دھیان کھِچدا اے۔ ایہہ
\VUgras{شکاری پرندہ}، اپݨی تکڑی \VUgras{چُنج} نال، اکثر اپݨے بچیاں لئی کوئی لیلا
یا بکری دا بچہ اپݨے \VUgras{پنجیاں} وچ دبائی لیاندا اے۔ اوہدے \VUgras{پر} اِنّے
وڈے نیں پئی اوہ اپݨے بھارے بوجھ نال دیر تیکر تر سکدا اے۔

%%K t09-tit-4 sub
\VUsaut{3.0mm}
\VUpk{10.2pt}{40}{پیدل ٹُرن والا۔}
\VUsaut{3.0mm}

%%K t09-c1-06-1 p
6. --- ایتھے دی سب توں سوہݨی سیر کو دے \VUgras{چرچرے} \textsuperscript{(10)} تیکر
جاݨا اے۔ \VUgras{ڈِگدا پاݨی} وَروَلا بݨ کے تھلے آندا اے تے فیر شہر دے لہندے دے
\VUgras{دیودار دے جنگل} \textsuperscript{(11)} وچ اک سوہݨے \VUgras{نالے} دے رُوپ وچ
گُم ہو جاندا اے۔ اسیں ایس جنگل وچوں \VUgras{موسم دے مرکز} \textsuperscript{(12)}
تیکر چڑھے، تے \VUgras{منارے} دے اُتوں اسیں علاقے دا سارا \VUgras{پنچھی جیہا منظر}
سمیٹ لیا۔ \VUgras{قدرتی منظر} بہت ودھیا اے، تے لگدا اے \VUgras{قدرت} نے ایس دیس
اُتے اپݨے سارے خزانے لُٹا دِتے نیں۔

%%K t09-c1-07-1 p
7. --- تھلے اُترن لئی اسیں \VUgras{رسے والی ریل} \textsuperscript{(13)} لئی تے اک
نِکے \VUgras{اسٹیشن} اُتے رُکے، اک پرانے \VUgras{جاگیردار قلعے} دے
\VUgras{کھنڈراں} \textsuperscript{(14)} کول، جیدے وچ کدے پراتز دے جاگیردار رہندے سن۔

%%K t09-c1-08-1 p
8. --- وچارا قلعہ! اوہ کیہ سوچدا ہووے گا جدوں اپݨے تھلے اوس سُتھرے
\VUgras{پاݨیاں والے شہر} \textsuperscript{(15)} نوں ویکھدا اے، جیہڑا اج اک فیشن دا
\VUgras{گرم چشمیاں دا صحت گھر} اے؟

%%K t09-c1-08-2 p
\VUgras{موسم} اِنّا سوہݨا اے تے \VUgras{کھݨج والا پاݨی} اِنّا گُݨ کاری پئی
\VUgras{صحت گھر} \textsuperscript{(17)} وچ لیا \VUgras{علاج دا دور} سچ مُچ سُکھ دی گل
اے۔

%%K t09-tit-5 sub
\VUsaut{3.0mm}
\VUpk{10.2pt}{40}{اک سوہݨا پاݨیاں والا شہر۔}
\VUsaut{3.0mm}

%%K t09-c1-09-1 p
9. --- سچ مُچ، ٹِکاء نوں سوہݨا بݨاؤݨ لئی سب کجھ کیتا گیا اے۔ ریل لے جاوݨ لئی اک
وڈا \VUgras{پُل} \textsuperscript{(16)} بݨایا گیا؛ \VUgras{پاݨیاں دے گھر} دا
\VUgras{باغ} \textsuperscript{(18)}، جِتھے \VUgras{موسیقی دیاں محفلاں} ہوندیاں نیں،
بڑے سلیقے نال سجایا گیا اے۔ کئی سوہݨیاں \VUgras{کوٹھیاں} \textsuperscript{(19)}
\VUgras{جوا گھر} \textsuperscript{(21)} دے چوہاں پاسیں بݨیاں نیں، تے اک بہت چنگی
طرحاں رکھی \VUgras{شاہ راہ} \textsuperscript{(22)} آوݨ جاوݨ بہت سوکھا کر دیندی اے۔

%%K t09-c1-10-1 p
10. --- اک سادہ \VUgras{بند} \textsuperscript{(23)} \VUgras{سڑک}
\textsuperscript{(20)} نوں اصلی \VUgras{وادی} \textsuperscript{(25)} توں وکھ کردا اے،
جیدے تلے وچ اک سوہݨی نِکی \VUgras{جھیل} \textsuperscript{(26)} اے، جیہڑی اک صاف
\VUgras{نالے} \textsuperscript{(24)} نوں پاݨی دیندی اے۔

%%K t09-c1-11-1 p
11. --- وادی دے دوجے پاسے لوہے دی اک \VUgras{کان} \textsuperscript{(28)} پُٹی جا
رہی اے۔ کئی واری میں \VUgras{کان کنّاں} نوں \VUgras{کھوہ} وچ اُتردے ویکھݨ گیا، تے
میں اوس \VUgras{سُرنگ} وچ وی گیا جِتھے اوہ سارا دن اوہ \VUgras{کچی دھات} کڈدے نیں
جیدے وچ \VUgras{دھات} رہندی اے۔

%%K t09-c1-12-1 p
12. --- کان کول اک وڈی \VUgras{ڈھلائی شالہ} بݨائی گئی، تاں جو کڈی ہوئی دولت اوتھے
ای کم آ جاوے۔ \VUgras{بھٹی} \textsuperscript{(29)}، جیہڑی ایتھے بہتات نال لبھدے
\VUgras{کوئلے} نال گرم کیتی جاندی اے، چھیتی تے سستے وچ \VUgras{کچا لوہا} لبھݨ
دیندی اے۔

%%K t09-c1-13-1 p
13. --- کان توں دور نئیں اک \VUgras{کھان} \textsuperscript{(30)} کٹی جا رہی اے۔
بُڈھا \VUgras{کھان دا مزدور} اپݨی \VUgras{کہی} نال جیہڑا توڑدا اے اوہ نہ
\VUgras{گرینائٹ} اے، نہ \VUgras{پورفری}، نہ \VUgras{ریتلا پتھر}، سگوں گھر بݨاؤݨ دا
سدھارن \VUgras{عمارتی پتھر}۔

%%K t09-c1-14-1 p
14. --- ایس دیس دیاں سب توں سوہݨیاں سوغاتاں وچوں اک اے \VUgras{دیودار}
\textsuperscript{(31)}۔ ہر تھاں — کِسے \VUgras{کھڈ} \textsuperscript{(32)} دے تلے وچ،
کِسے \VUgras{ابشار} \textsuperscript{(33)} دے کنڈھے، کِسے \VUgras{غار}
\textsuperscript{(34)} یا کگر کول — اوہدیاں پتلیاں سُوئیاں اپݨا ہرا سُر جوڑ
دیندیاں نیں۔

%%K t09-tit-6 sub
\VUsaut{3.0mm}
\VUpk{10.2pt}{40}{پیدل ٹُرنا۔}
\VUsaut{3.0mm}

%%K t09-c1-15-1 p
15. --- میں اپݨیاں چھٹیاں دا فائدہ لمیاں سیراں لئی چُک رہیا واں۔ جیہڑیاں
\VUgras{پگڈنڈیاں} \textsuperscript{(35)} پہاڑ اُتے چکر کٹدیاں چڑھدیاں نیں، اوہناں
اُتے کئی \VUgras{کلومیٹر}، تے اکثر کئی \VUgras{کوہ}، پندھ پتہ لگے بغیر مُک جاندا
اے۔

%%K t09-c1-15-2 p
\VUgras{میل پتھر} \textsuperscript{(36)} تے \VUgras{سیدھ سُوچک} \textsuperscript{(37)}
راہواں نوں نشان لاندے نیں، جیہناں اُتے اکثر کوئی جوان \VUgras{پہاڑی مُنڈا}
\textsuperscript{(38)} ملدا اے، بغل وچ \VUgras{جھولا} \textsuperscript{(40)} چُکی تے اک
\VUgras{مارموٹ} \textsuperscript{(39)} چُکیا، یا کوئی \VUgras{بنجارا}
\textsuperscript{(41)}، جیدی \VUgras{جھلّ والی ٹوپی} \textsuperscript{(42)} اوہدے
پھٹے پرانے \VUgras{اوڑھݨ} \textsuperscript{(43)} نال عجیب میل کھاندی اے۔ اوہ اک
سِدھائے \VUgras{رِچھ} \textsuperscript{(44)} نوں \VUgras{سنگل} نال لئی ٹُردا اے۔

%%K t09-tit-7 sub
\VUpk{10.2pt}{40}{چڑھائی۔}
\VUsaut{3.0mm}

%%K t09-c1-16-1 p
16. --- تقریباً ہر دن میں اک \VUgras{گائیڈ} \textsuperscript{(48)} نال
\VUgras{چڑھائی} کرن جانا واں، تے مینوں بڑا مان ہوندا اے جدوں پِٹھ اُتے اپݨا
\VUgras{پٹھو} \textsuperscript{(47)} تے ہتھ وچ اپݨی \VUgras{لاٹھی} چُکی، کِسے سچے
\VUgras{سیلانی} \textsuperscript{(46)} وانگوں، میں \VUgras{چٹاناں}
\textsuperscript{(45)} اُتے چڑھائی کرنا واں۔

%%K t09-c1-17-1 p
17. --- پہاڑ دی چوٹی توں میدان دے لوک کِیڑیاں توں وڈے نئیں لگدے؛
\VUgras{بھیڈاں دا اِجّڑ} \textsuperscript{(50)}، جیہڑا \VUgras{چراگاہ}
\textsuperscript{(49)} وچ چردا اے، نِکیاں دا کھڈوݨا لگدا اے۔ اک \VUgras{چیڑ}
\textsuperscript{(51)}، یا \VUgras{لکڑ دا پہاڑی گھر} \textsuperscript{(52)} تیکر، ہرے
اُتے اک بندے توں ودھ نئیں۔

%%K t09-c1-18-1 p
18. --- ایہناں سیراں وچ سانوں \VUgras{پگڈنڈی} \textsuperscript{(53)} اُتے اکثر کوئی
\VUgras{شاموآ دا شکاری} \textsuperscript{(54)} ملدا اے، جیہڑا ہُݨے ہُݨے ماریا ہویا
ڈنگر مونڈھے اُتے چُکی ہوندا اے۔

%%K t09-c1-19-1 p
19. --- میں اپݨے بُوٹیاں دے سنگرہ وچ \VUgras{فرن} \textsuperscript{(55)} دا اک بہت
انوکھا پتّر تے \VUgras{ہیدر} \textsuperscript{(57)} دی اک سوہݨی گلابی ٹہݨی رکھی اے،
جیہڑی میں پچھلی سیر وچ اک نِکی \VUgras{غار} \textsuperscript{(56)} دے مونہہ اُتے
کٹھی کیتی سی۔

%%K t09-tit-8 sub
\VUsaut{3.0mm}
\VUcentre{12.0pt}{0}{\textit{دوجا منظر۔}}
\VUsaut{3.0mm}

%%K t09-tit-9 sub
\VUpk{11.4pt}{40}{جنگل۔ --- شکار۔}
\VUsaut{3.0mm}

%%K t09-c2-01-1 p
1. --- کل میں نے جنگل وچ اک آنند بھریا دن لنگھایا۔ اوتھے رہݨ والے ڈنگراں تے
\VUgras{کِیڑیاں} \textsuperscript{(5)} دے رُجھے جیوݨ نے مینوں بہت چنگا لگا۔

%%K t09-c2-01-2 p
\VUgras{مکڑی} \textsuperscript{(1)}، جیہڑی دو ٹہݨیاں وچکار اپݨا \VUgras{جالا}
\textsuperscript{(2)} بُݨ کے لنگھدی \VUgras{مکھی} \textsuperscript{(3)} پھڑدی اے،
\VUgras{گھونگا} \textsuperscript{(4)}، جیہڑا اپݨا کھول چُکی اوکھا ای اگے ودھدا اے،
اوہ بھونرا جیہڑا اپݨے شکار دی ٹوہ وچ اے، اوہ محنتی کِیڑی جیہڑی \VUgras{بھون}
\textsuperscript{(6)} کول کم وچ جُٹی اے — ایہہ سارے نِکے جیو اوکھے جوش نال آندے
جاندے تے کم کردے رہے۔

%%K t09-c2-02-1 p
2. --- اک \VUgras{سپ} \textsuperscript{(7)}، جیہنوں پہلاں میں \VUgras{ناگ} سمجھیا
پر جیہڑا نِرا اک بے قصور \VUgras{پاݨی دا سپ} سی، اک مُڈھ دے تھلے کُنڈلی مارے سی،
تے وِچ وِچ اپݨا پتلا سر کڈدا سی، جیہڑا ہمیشہ ڈنگݨ نوں تیار لگدا سی۔ میرے سامݨے
اک وڈا \VUgras{سلگ} \textsuperscript{(8)} بھاری چال نال رِڑھ رہیا سی، اوہناں سوادی
نِکیاں \VUgras{جنگلی اسٹرابیریاں} \textsuperscript{(11)} وچوں اک چٹ کر لیݨ مگروں،
جیہڑیاں ایتھے بہتات نال اُگدیاں نیں۔

%%K t09-c2-03-1 p
3. --- زمین \VUgras{جنگلی پھُلاں} نال وِچھی سی : \VUgras{پرم روز}
\textsuperscript{(9)}، \VUgras{پیری ونکل} \textsuperscript{(10)} تے ہور وی بہت سارے
بُوٹے، جیہناں نوں میں جاݨدا نہ سی، \VUgras{گھاہ دے گُچھیاں} \textsuperscript{(12)}
وچکار کھِڑے سن۔

%%K t09-c2-04-1 p
4. --- ایس علاقے وچ \VUgras{ٹرفل} تقریباً اوہنی ای عام اے جِنّی \VUgras{کھمبی}
\textsuperscript{(14)}، تے اک جنگل دے رکھوالے دا \VUgras{سُور دا بچہ}
\textsuperscript{(13)} لالچ نال لبھݨ وچ لگا سی پئی کجھ لبھ جاوے۔

%%K t09-tit-10 sub
\VUsaut{3.0mm}
\VUpk{10.2pt}{40}{چوری دا شکار۔}
\VUsaut{3.0mm}

%%K t09-c2-05-1 p
5. --- میں نے اک منظر دا \VUgras{انت} ویکھیا جیہنے مینوں بہت ہسایا۔ اپݨے
\VUgras{کھوجی کُتے} \textsuperscript{(15)} دے نک دی بدولت \VUgras{جنگل دے رکھوالے}
\textsuperscript{(16)} نے ہُݨے ہُݨے اک \VUgras{چور شکاری} \textsuperscript{(22)} نوں
اوس ویلے رنگے ہتھیں پھڑیا سی جدوں اوہ اپݨا ماریا \VUgras{شکار}
\textsuperscript{(17)} چُک لے جاوݨ دی تیاری وچ سی۔ پھڑے جاوݨ نالوں چنگا اوہنوں اپݨا
شکار چھڈ دیݨا لگا، سو چور شکاری نس گیا، تے \VUgras{سہا} \textsuperscript{(18)}،
\VUgras{ہرنی} \textsuperscript{(19)}، \VUgras{بارہ سنگھا} \textsuperscript{(20)} تے
\VUgras{جنگلی سُور} \textsuperscript{(21)} گھاہ اُتے پئے رہ گئے، جنگل دے رکھوالے دی
خوشی لئی۔

%%K t09-c2-06-1 p
6. --- جنگل وچ رُکھاں دیاں جِنّیاں قسماں نیں، اوہ دنگ کر دیندیاں نیں۔
\VUgras{بیچ} \textsuperscript{(23)} تے \VUgras{بھوج پتّر} \textsuperscript{(26)}،
\VUgras{چیڑ}، جیہناں توں \VUgras{گوند} \textsuperscript{(27)} لبھدی اے،
\VUgras{بلوط} \textsuperscript{(29)} تے ہور وی بہت سارے اپݨیاں ٹہݨیاں ملاندے نیں۔

%%K t09-c2-06-2 p
اُچے رُکھاں دے مُڈھاں نال لپٹی \VUgras{آئیوی} \textsuperscript{(24)}
\VUgras{اُچے جنگل} \textsuperscript{(25)} نوں اک چمکدا ہرا رنگ دیندی اے، جیہڑا اوس
\VUgras{کائی} \textsuperscript{(31)} دی فِکی کُونی چھاں نال سوہݨے ڈھنگ نال مل جاندا
اے جیدے وچ \VUgras{ہرا کٹھ پھوڑا} \textsuperscript{(30)} کِیڑے لبھدا اے۔

%%K t09-tit-11 sub
\VUsaut{3.0mm}
\VUpk{10.2pt}{40}{جنگل دا منظر۔}
\VUsaut{3.0mm}

%%K t09-c2-07-1 p
7. --- پرانے بلوط مینوں موہ لیندے نیں۔ اوہناں دیاں وڈیاں وَل کھاندیاں
\VUgras{جڑاں} \textsuperscript{(32)}، اَدھ زمین توں باہر، اوہناں نوں زور تے تیج دا
شاندار رُوپ دیندیاں نیں، تے مینوں اچرج ہوندا اے پئی \VUgras{مالک}
\textsuperscript{(35)} اوہناں نوں وڈھوا کے اوس \VUgras{آری} \textsuperscript{(34)} نوں
سونپݨ اُتے کِویں راضی ہو جاندا اے جیہڑی \VUgras{لکڑہارے} \textsuperscript{(33)} دی
اے۔ وچارے \VUgras{مُڈھ} \textsuperscript{(36)}، اپݨی \VUgras{چھِلّ}
\textsuperscript{(37)} توں چھِلے یا \VUgras{کہاڑے} \textsuperscript{(38)} نال چیرے،
جیہڑا \VUgras{دیہاڑی دار مزدور} \textsuperscript{(39)} دا اے، مینوں دُکھ دیندے نیں۔

%%K t09-c2-08-1 p
8. --- کول ای اک بُڈھے \VUgras{کوئلہ بݨاؤݨ والے} \textsuperscript{(41)} نے اپݨی
\VUgras{کہی} نال کوئلے دا اک \VUgras{ڈھیر} \textsuperscript{(42)} لایا ہویا سی، تے اک
بُڈھی وڈھے رُکھاں دیاں ٹہݨیاں توں بݨی \VUgras{سُکی لکڑ} دی اک \VUgras{گٹھڑی}
\textsuperscript{(40)} لے جا رہی سی۔

%%K t09-tit-12 sub
\VUsaut{3.0mm}
\VUpk{10.2pt}{40}{شکار۔}
\VUsaut{3.0mm}

%%K t09-c2-09-1 p
9. --- میں \VUgras{جنگل دے رکھوالے دے گھر} \textsuperscript{(46)} وچ تھوڑا ساہ لیݨ
جاوݨ والا سی، پر اوس نِکے \VUgras{تال} \textsuperscript{(43)} تیکر اپڑ کے، جیدے اُتے
اک سوہݨا \VUgras{ہنس} \textsuperscript{(45)} تر رہیا سی، میں ویکھیا پئی ہاݨ دے اک
\VUgras{ہڑھ} \textsuperscript{(44)} نے راہ نوں سچ مُچ دا \VUgras{چِکڑ} بݨا دِتا اے۔
مینوں چکر کٹݨا پیا، تے \VUgras{دیودار} \textsuperscript{(49)}، \VUgras{سفیدیاں}
\textsuperscript{(48)} تے \VUgras{ایلم} \textsuperscript{(47)} کولوں لنگھدے ہوئے،
جیہڑے جنگل دے کنڈھے کھڑے نیں، میں ویکھیا پئی اک \VUgras{جھاڑی}
\textsuperscript{(54)} وچوں نکل کے اک شاندار \VUgras{بارہ سنگھا} \textsuperscript{(50)}
نسیا، جیدے پچھے اک نہ تھکݨ والا \VUgras{شکاری کُتیاں دا جھُنڈ} \textsuperscript{(51)}
لگا سی، جیہنوں \VUgras{نرسنگھا} \textsuperscript{(53)} اُکسا رہیا سی،
\VUgras{گھوڑیاں اُتے سوار شکاریاں} \textsuperscript{(52)} دا۔ وچارا ڈنگر بالکل تھک
چُکیا سی۔ اوہ میرے توں کجھ ای قدماں اُتے مردا ہویا ڈِگ پیا۔

%%K t09-c2-10-1 p
10. --- جنگل دے رکھوالے دا پُتر، جیہنوں \VUgras{شکار} دے رولے نے باہر کھِچ لیا، گھروں
نکلیا، تے اسیں دوویں مُڑ جنگل وچ آئے۔ اوہنے اک پرانے بلوط دی ٹہݨی اُتے اک
\VUgras{نیل کنٹھ} \textsuperscript{(56)} ویکھیا سی۔ اوہنوں لگا پئی اوہدا آہلݨا
\VUgras{پتّراں} \textsuperscript{(58)} وچ لُکیا ہووے گا، تے اوہ اوہنوں لیݨا چاہندا سی۔

%%K t09-c2-10-2 p
\VUgras{گلہری} \textsuperscript{(61)} وانگوں پھرتیلا، اوہ اک \VUgras{ٹہݨی}
\textsuperscript{(55)} توں دوجی ٹہݨی چڑھدا گیا، پر اوہنوں کجھ نہ لبھیا تے اوہ نِرا اک
بُڈھے \VUgras{کاں} \textsuperscript{(60)} نوں چونکا سکیا، جیہڑا اک \VUgras{شاخ}
\textsuperscript{(57)} اُتے بیٹھا سی۔
\VUsaut{4.0mm}
\VUfilet{22mm}
